\section{Length of curves and Riemann distance}

Let $M$ be a Riemannian manifold, $I \subset \mathbb{R}$ be an interval, and $c:I\to M$ a $C^{\infty}$ curve.
\begin{definition}[Length of a curve]\label{def:length_of_curve}
For $a,b\in I$ such that $a<$, the \emph{length} $L(c|_{[a,b]})$ of $c$ on $[a,b]$ is defined by
\[
L(c|_{[a,b]}) := \int^b_a ||c'(t)|| dt.
\]
\textcolor{red}{Here, I believe $||\dots|| = \sqrt{\langle \dots, \dots \rangle}$.}
\end{definition}

\begin{definition}[Regular curve]
A $C^{\infty}$ curve $C:I\to M$ is \emph{regular} if
\[
c'(t) \neq 0, \quad \forall t \in I.
\]
\end{definition}

\begin{theorem}[Inverse function theorem for regular curves] \label{thm:inverse_function_theorem}
Assume that the parameter $t$ is a $C^{\infty}$ function of $u$ on $J \subset \mathbb{R}$, i.e. $t = t(u)$ satisfying:
\begin{enumerate}
\item $t: J \to I$ is surjective.
\item $\frac{dt}{du}(u) \neq 0 \forall u \in J$.
\end{enumerate}
Then there exists a $C^{\infty}$ inverse $t^{-1}: I \to J$. This is called a \emph{transformation of parameter} $t$ to $u$.
\end{theorem}

Let $\tilde{c} := c(t(u))$ be a reparametrization of $c$ by $t$. Fix a chart $(U, \phi), \phi = (x^1, \dots, x^n)$ and subinterval $I'\subset I$ such that $c(I')\subset U$.
For $t \in I'$ and $u\in t^{-1}(I')$, let
\[
\phi(c(t)) = (c^1(t), \dots, c^n(t)), \quad \phi(\tilde{c}(u)) = (\tilde{c}^1(u), \dots, \tilde{c}^n(u)).
\]
Since $\tilde{c}^i(u) = c^i(t(u))$, we have
\[
\begin{aligned}
\tilde{c}'(u) &= \sum_{i=1}^n \frac{d\tilde{c}^i}{du}\left(\frac{\partial}{\partial x^i}\right)_{\tilde{c}(u)} \\
&= \sum_{i=1}^n \frac{dt}{du} \frac{dc^i}{dt}(t(u))\left(\frac{\partial}{\partial x^i}\right)_{\tilde{c}(u)}\\
&= \frac{dt}{du} c'(t(u)).
\end{aligned}
\]
Since $(U,\phi)$ was arbitrary, for all $u \in J$,
\[
\tilde{c}'(u) = c'(t(u))\frac{dt}{du}.
\]
\begin{proposition}
Let $c:I\to M$ be a regular curve and $t = t(u)$ ($u \in J$) be a transformation. Then
\begin{enumerate}[label=(\arabic*)]
\item If $\frac{dt}{du}(u) > 0 $ ($\forall u \in J$), then $L(\tilde{c}|_{[a,b]}) = L(c|_{[t(a),t(b)]}) \forall a,b \in J (a < b)$.
\item If $\frac{dt}{du}(u) < 0 $ ($\forall u \in J$), then $L(\tilde{c}|_{[a,b]}) = L(c|_{[t(b),t(a)]}) \forall a,b \in J (a < b)$.
\end{enumerate}
\end{proposition}
\begin{proof}
\begin{enumerate}[label=(\arabic*)]
\item By \autoref{def:length_of_curve} (length of a curve),
\[
\begin{aligned}
L(\tilde{c}|_{[a,b]}) &= \int^b_a ||\tilde{c}'(u)|| du \\
&= \int^b_a ||c'(t(u))|| \left|\frac{dt}{du}\right| du \\
&= \int^{t(b)}_{t(a)} ||c'(t)|| dt \\
&= L(c|_{[t(a),t(b)]}).
\end{aligned}
\]
\item Again, by \autoref{def:length_of_curve},
\[
\begin{aligned}
L(\tilde{c}|_{[a,b]}) &= \int^b_a ||\tilde{c}'(u)|| du \\
&= \int^b_a ||c'(t(u))|| \left|\frac{dt}{du}\right| du \\
&= \int^{t(a)}_{t(b)} ||c'(t)|| dt \\
&= L(c|_{[t(b),t(a)]}).
\end{aligned}
\]
\end{enumerate}
\end{proof}

\begin{definition}[Arclength parameter]
A regular curve $c:I\to M$ is said to have an \emph{arclength parameter} $s$ if $||c'(s)|| = 1$ for all $s \in I$.
\end{definition}
If $c:I\to M$ has an arclength parameter, then
\[
L(c|_{[a,b]}) = \int^b_a ||c'(s)|| ds = \int^b_a 1 ds = b-a, \quad \forall a,b \in I (a < b).
\]
Meaning the "elapsed time" $b-a$ equals the distance traveled along $L(c|_{[a,b]})$.
For any regular curve $c:I\to M$ and $t_0 \in I$, let
\[
s(t) := \int^t_{t_0} ||c'(t)|| dt.
\]
Since $\frac{ds}{dt} = ||c'(t)|| > 0 \implies$ there exists a global inverse $t=t(s)$ by \autoref{thm:inverse_function_theorem} (inverse function theorem).

\begin{lemma}
$\tilde{c}(s) = c(t(s))$ has an arclength parameter $s$.
\end{lemma}
\begin{proof}
By $\tilde{c}'(s) = c'(t(s)) \frac{dt}{ds}$, we have
\[
\begin{aligned}
||\tilde{c}'(s)|| &= ||c'(t(s))|| \left|\frac{dt}{ds}\right|,\\
\left|\frac{dt}{ds}\right| &= \frac{1}{\frac{ds}{dt}} = \frac{1}{||c'(t)||}.
\end{aligned}
\]
Meaning we have $||\tilde{c}'(s)|| = 1$ for all $s \in I$.
\end{proof}

\begin{definition}[Piecewise smooth curve]
A continuous curve $c:[a,b]\to M$ is a \emph{piecewise smooth regular curve} if there exists a finite sequence $a = t_0 < t_1 < \dots < t_k = b$ such that
\[
c|_{[t_,t_{i+1}]}:[t_i,t_{i+1}] \to M
\]
is regular for all $i = 0, \dots, k-1$.

Furthermore, the length $L(c)$ is defined by
\[
L(c) := \sum^{k-1}_{i=0} L(c|_{[t_i,t_{i+1}]}) = \sum^{k-1}_{i=0} \int^{t_{i+1}}_{t_i} ||c'(t)|| dt.
\]
\end{definition}

Assume $M$ is connected, i.e. any two points are joinable by a regular curve.
\begin{definition}[Riemannian distance]
The \emph{Riemannian distance} $d(p,q)$ between two points $p, q \in M$ is defined as
\[
d(p,q) := \inf \{ L(c) : c \text{ is a piecewise smooth regular curve from } p \text{ to } q \}.
\]
$d:M\times M \to \mathbb{R}$ is called the \emph{Riemannian distance function}.
If there exists some function $c$ attaining the "inf", then $c$ is called a \emph{minimal curve}.
\end{definition}

\begin{theorem}
The function $d:M\times M \to \mathbb{R}$ is satisfies the axioms of distance. That is, for all $p,q,r \in M$,
\begin{enumerate}[label=(\arabic*)]
\item $d(p,q) = 0$ if and only if $p=q$,
\item $d(p,q) = d(q,p)$,
\item $d(p,r) \leq d(p,q) + d(q,r)$.
\end{enumerate}
\label{thm:distance_axioms}
\end{theorem}

\begin{lemma}
$\left(U, \phi \right)$ : a chart, $D \subset \phi(U)$ a cpt set.
$c : [a,b] \rightarrow \phi^{-1} (D)$ : a piecewise smooth regular curve, Then $\exists C = C\left( (U, \phi), g_{ij}, D \right) >0$ s.t 
for $\bar{c} := \phi \circ c :[a,b] \rightarrow \mathbb{R}^n$
\[
\begin{aligned}
    C^{-1} L(\bar{c}) &\leq L(c) \leq C L(\bar{c}), \text{ where}\\
    L(\bar{c}) &= \int_a^b \|\bar{c}'(t)\|_{\mathbb{R}^n} dt \text{ is the Euclidean length.}
\end{aligned}
\]
\label{lemma:length_of_curve_in_chart}
\end{lemma}

\begin{proof}
    Define a continous function $f: \phi(u) \times \mathbb{R}^n \rightarrow \mathbb{R}$ by
    \[
    \begin{aligned}
        f(x,v) &:= \sum_{i,j = 1}^n v^i v^j g_{ij} (\phi^{-1}(x))
    \end{aligned}
    \]
    Since $D\times \mathbb{S}^{n-1} \subset \mathbb{R}~{2n}$ is cpt, $f|_{D \times \mathbb{S}^{n-1}}$ has max/min. Since $f>0$, the min is positive.
    \[
    \begin{aligned}
        \implies \exists C>0 \text{ s.t } c^{-2} \leq f \leq C^2 \text{ on } D \times \mathbb{S}^{n-1}
    \end{aligned}
    \]
    Thus, $\forall x \in D$ and $\forall v \in \mathbb{R}^n \setminus \{0\}:$
    \[
    \begin{aligned}
        C^{-2} \leq f\left(x, \frac{v}{\|v\|_{\mathbb{R}^n}} \right) \leq C^2
    \end{aligned}
    \]
    Multiplying $\|v\|^2_{\mathbb{R}^n}$
    \[
    \begin{aligned}
        C^{-2} \|v\|_{\mathbb{R}^n}^2 \leq \sum^n_{i,j=1} v^i v^j g_{ij}(\phi^{-1}(x)) \leq C^2 \|v\|_{\mathbb{R}^n}^2
    \end{aligned}
    \]
    Let $\bar{c} = \left(c^1(t), \dots, c^n(t)\right) \implies c'(t) = \sum^n_{i=1} \frac{dc^i}{dt} \partial_i$.
    \[
    \begin{aligned}
        \bar{c}'(t) &= \left(\frac{dc^1}{dt}, \dots, \frac{dc^n}{dt} \right). \text{ Set } v := \bar{c}'(t) \\
        &C^{-2} \| \bar{c}'(t)\|_{\mathbb{R}^n}^2 \leq \|c'(t)\|^2 \leq C^2 \|\bar{c}'(t)\|_{\mathbb{R}^n}^2
    \end{aligned}
    \]
    Taking square roots and integrating with respect to $t$ yields the result.
\end{proof}

\begin{lemma}
    Let $(U, \phi)$ be a chart of M.
    Choose $x_0 \in \phi(U)$ and $r>0$ s.t $D := \bar{B}(x_0, 3r) \subset \phi(U)$.
    Let $C>0$ be the constant from \autoref{lemma:length_of_curve_in_chart}. Then:
    \begin{enumerate}
        \item $\forall p,q \in \phi^{-1}({B}(x_0, r))$: \\
        $C^{-1} \| \phi(p) - \phi(q) \|_{\mathbb{R}^n} \leq d(p,q) \leq C \|\phi(p) - \phi(q)\|_{\mathbb{R}^n}$.
        \item $\forall p \in \phi^{-1}({B}(x_0, r)), \forall q \in M \setminus \phi^{-1}(D) : d(p,q) \leq 2C^{-1}r$
    \end{enumerate}
    \label{lemma:distance_in_chart}
\end{lemma}

\begin{proof}
    (1) $p,q \in \phi^{-1}({B}(x_0, r))$. $\bar{\gamma}$ : the line segment in $\mathbb{R}^n$ connecting $\phi(p)$ and $\phi(q)$. 
    Since $B(x_0, r)$ is convex, $\bar{gamma} \subset B(x_0, r)$. Set $\gamma := \phi^{-1} \circ \bar{\gamma}$.
    By \autoref{lemma:length_of_curve_in_chart},
    \[
    \begin{aligned}
        \|\phi(p) - \phi(q) \|_{\mathbb{R}^n} \leq L(\bar{\gamma}) \leq C^{-1} L(\gamma) \geq C^{-1} d(p,q).
    \end{aligned}
    \]
    Next, take any piecewise smooth refular curve $C: [0,l] \to M$ from $p$ to $q$. 
    If $c([0,l]) \subset \phi^{-1}(D)$, let $\bar{c} := \phi \circ c$. By \autoref{lemma:length_of_curve_in_chart},
    \[
    \begin{aligned}
        L(c) \geq C^{-1} L(\bar{c}) \geq C^{-1} \|\phi(p) - \phi(q)\|_{\mathbb{R}^n}.
    \end{aligned}
    \]
    If $c([0,l]) \not\subset \phi^{-1}(D)$, let
    \[
    \begin{aligned}
        t_1 := sup \{t\in [0,l] | c([0,t]) \subset \phi^{-1}(D)\}
    \end{aligned}
    \]
    and $c_1 := c|_{[0,t_1]}$. By \autoref{lemma:disjoint_or_equal}, $c_1(t_1) \in \phi^{-1}(\partial D)$
    Let $\bar{c}_1 := \phi \circ c_1$. By \autoref{lemma:length_of_curve_in_chart},
    \[
    \begin{aligned}
        L(c) \geq L(c_1) \geq C^{-1} L(\bar{c}_1).
    \end{aligned}
    \]
    Since $\phi(p), \phi(q) \in B(x_0, r )$ and $\bar{c_1}(t_1) \in \partial B(x_0, 3r)$
    \[
    \begin{aligned}
        &L(\bar{c_1}) \geq 2r > \|\phi(p)-\phi(q)\|_{\mathbb{R}^n} \\
       \implies &L(c) > C^{-1} \|\phi(p) - \phi(q)\|_{\mathbb{R}^n} \text{ in both cases}
    \end{aligned}\]
    Taking inf over c (using the fact that C in \autoref{lemma:length_of_curve_in_chart} is independent of c)
    \[
    \begin{aligned}
        d(p,q) \geq C^{-1} \|\phi(p) - \phi(q) \|_{\mathbb{R}^n}
    \end{aligned}
    \]
    (2) Prf is similar to (1) above (omitted).
\end{proof}

\begin{proof}[Proof of \autoref{thm:distance_axioms}]
    (1) $p=q \implies$ taking $c(t) = p$ gives $L(c) = 0 $ thus $d(p,q) =0$.
    conversely, assume $d(p,q) = 0$.
    Fix a xhart $(U, \phi)$ around $p$ and let $x_0 := \phi(p)$. 
    Choose $r>0$ s.t $D_r := \bar{B}(x_0, 3r) \subset \phi(U)$. By \autoref{lemma:distance_in_chart}, 
    \[
    \begin{aligned}
        \exists C_r > 0 \text{ s.t if } q \in M \setminus \phi^{-1} (D_r), \text{ then } d(p,q) \geq 2C_r^{-1} r
    \end{aligned}
    \]
    Contradicting $d(p,q) = 0$. Thus $q \in \phi^{-1}(D_r)$.
    Since $r>0$ is arbitrary, $p=q$

    (2) Let $c$ be a piecewise smooth regular curve from $p$ to $q$. Reversing, the parameter yiels a curve from $q$ to $p$ with the same length. $\implies d(p,q) = d(q,p)$

    (3) Let $c_{pq}$ be a curve from $p$ to $q$, and $c_qr$ from $q$ to $r$. Concatenating them yields a piecewise smooth regular curve from $p$ to $r$.
    $d(p,r) \leq L(c_{pqr}) = L(c_{pq}) + L(c_{qr})$ Taking inf over $c_{pq}$ and $c_{qr} \implies$ triangle inequality. 
\end{proof}

\begin{remark}
    Hausdorff property is essential for $d(p,q) = 0$. In non-Hausdorf manifolds, $d(p,q)=0$ does not necessarily imply $p=q$
\end{remark}

\begin{problem}
    Define $L, R_+, R_- \subset \mathbb{R}^2$ as:
    \[
    \begin{aligned}
        &L:= \{(x, 0) \in \mathbb{R}^2 | x<0\}\\
        &R_+ := \{(x, 1) \in \mathbb{R}^2 | x \geq 0\}, R_- := \{(x, -1) \in \mathbb{R}^2 | x \geq 0\}
    \end{aligned}
    \]
    $M:=L \cup R_+ \cup R_-$, $U:= L \cup R_+, V:= L \cup R_-$
    Let $\pi_1 : \mathbb{R}^2 \to \mathbb{R}$ be the projectioon to the 1st component
    and $\phi := \pi_1 |_U : U \to \mathbb{R}, \psi := \pi_1 |_V : V \to \mathbb{R}$
    $W \subset M$ is open : $\equiv \phi(W \cap U), \psi(W \cap  V)\subset \mathbb{R}$ are open.
    This makes $M$ a topological space s.t. $\phi, \psi$ are homeomorphisms,
    \[
    \begin{aligned}
        \phi(U \cap V) = \psi(U\cap V) = (-\infty, 0), \text{ and } \psi \circ \phi^{-1} = Id_{(-\infty, 0)}
    \end{aligned}
    \]
    \begin{enumerate}
        \item Show M is not Hausdorff
        \item Show $d \left((0,1), (0,-1)\right) = 0$ for any riemann metric on $M$
    \end{enumerate}
\end{problem}

\begin{definition}
    For $p\in M$, $r>0$
    \[
    \begin{aligned}
        B(p,r) := \{q \in M | d(p,q) < r\}
    \end{aligned}
    \]
    is called a geodesic ball of radius $r$ centered at $p$
\end{definition}

\begin{proposition}
    For a subset $O \subset M$ the following are equivalent:
    \begin{enumerate}
        \item $O$ is open in $M$
        \item $\forall p \in O, \exists \delta > 0 \text{ s.t. } B(p, \delta) \subset O$
    \end{enumerate}
    \label{prop:open_in_manifold}
\end{proposition}

\begin{proof}
    Take a chart $(U, \phi)$ around $p$ and let $x_0 := \phi(p)$. Since $\phi(O \cap U) \subset \mathbb{R}^n$ is open,
    $\exists \gamma > 0$ s.t. $D_i = \bar{B}(x_0, 3r) \subset \phi(O \cap U)$. For the constant $C>0$ from \autoref{lemma:distance_in_chart}, choose $\delta >0$ s.t.
    $0 < \delta < 2 C^{-1}r$. Then $\forall q \in M \setminus \phi^{-1}(D),$
    $d(p,q) \geq 2C^{-1}r > \delta$. Thus $q \in B(p, \delta) \implies q \in \phi^{-1}(D)$,
    we have $B(p,\delta) \subset O$.
    $(2) \implies (1)$ : Let ${(U_{\alpha}, \phi_{\alpha})}_{\alpha \in \Lambda}$ be an atlas of $M$. Then $O = \bigcup_{\alpha \in \Lambda} (O \cap U_{\alpha})$.
    Since each $\phi_{\alpha} : U_{\alpha} \to \phi_{\alpha} (U_{\alpha})$ is a homeo,
    $O \cap U_{\alpha}$ is open in $U_{\alpha} \equiv \phi_{\alpha} (O \cap U_{\alpha})$ is open in $\phi_{\alpha}(U_{\alpha})$.
    Thus it suffices to show: $\phi (O \cap U) \subset \mathbb{R}^n $ is open for any chart $(U, \phi)$

    $x \in \phi(O \cap U)$, set $p := \phi^{-1}(x) \in O \cap U$. By assumption, $\exists r_1 >0 \text{ s.t. } B(p, r_1) \subset O$.
    Since $\phi (U) \subset \mathbb{R}^n$ is open, $\exists r_2>0 \text{ s.t. } \bar{B}(x, 3 r_2) \subset \phi(U)$.
    For $(U, \phi)$ and $D := \bar{B}(x, 3 r_2)$ let $C>0$ be the constant from \autoref{lemma:length_of_curve_in_chart}. Set $r:= min{C^{-1}r_1, r_2}$.
    Then $B(p, C r) \subset B(p, r_1) \subset O$. For any $y \in B(x, r)$ set $q := \phi^{-1}(y) \in \phi^{-1}(B(x. r_2)) \subset U$. By \autoref{lemma:length_of_curve_in_chart}:
    \[
    \begin{aligned}
        d(p,q) \leq C \| x-y \|_{\mathbb{R}^n} < C r,
    \end{aligned}
    \]
    Thus $q \in B(p, Cr) \subset O$, which implies $q \in O \cap U$ and $y \in \phi(O \cap U)$. Therefore $B(x,r) \subset \phi(O \cap U)$ and $\phi(O \cap U)$ is open in $\mathbb{R}^n$
\end{proof}

\begin{problem}
    Using \autoref{prop:open_in_manifold}, show that for any $p \in M$ and $r>0$, the geodesic ball $B(p,r)$ is an open set in $M$.
    The following properties hold for distance functions on general metric spaces, thus proofs are omitted. $d: M \times M \to \mathbb{R}$.
    \begin{proposition}
        $\forall p_0 \in M$, $f(p) := d(p, p_0)$ is continous
    \end{proposition}
    \begin{proposition}
        For a sequence ${p_i}_{i=1}^{\infty} \subset M$ and $p\in M$, the following are equivalent:
        \begin{enumerate}
            \item $p_i \to p$ as $i \to \infty$
            \item $d(p_i, p) \to 0$ as $i \to \infty$
        \end{enumerate}
        Basic topological properties of geodesic balls (proofs omitted:)
    \end{proposition}
    \begin{proposition}
        $\forall p \in M$ and $r>0$
        \begin{enumerate}
            \item $\partial B(p,r) = {q \in M | d(p,q)=r}$
            \item $\bar{B}(p,r) := \bar(B(p,r)) = {q \in M | d(p,r) \leq r}$
        \end{enumerate}
    \end{proposition}
\end{problem}